% !TeX root = ../tg.tex
\chapter{结论与展望}
\section{全文总结}
随着人工智能技术与移动互联网应用的深度融合，分布式学习范式已成为挖掘海量私有数据价值的关键途径。然而，在实际部署过程中，分割 - 联邦学习（Split Federated Learning, SFL）面临着数据质量低下与系统资源受限的双重挑战。低质量数据表现为非独立同分布（Non-IID）特性与标签噪声的共存，这严重影响了模型的收敛性与泛化能力；而资源受限则体现在边缘设备计算算力弱、内存容量小及通信带宽窄，限制了复杂深度学习模型的有效落地。针对上述问题，本文立足于数据与系统双重维度，围绕 SFL 在复杂现实场景中的高效、稳健应用开展了深入的理论与算法研究。本文的主要研究工作总结如下：
首先，在理论分析层面，本文填补了 SFL 在复杂低质量数据场景下收敛性理论支撑的空白。针对现有研究缺乏对标签噪声与 Non-IID 数据耦合影响量化的问题，本文构建了包含标签噪声的客户端数据集模型，并克服了 SFL 模型切分与更新过程复杂性带来的分析难点。本文首次推导了 SFL 在 Non-IID 数据和标签噪声共存条件下的收敛界，分别针对回归任务常用的均方误差（MSE）损失函数与分类任务常用的交叉熵（Cross-Entropy）损失函数进行了严谨的数学证明。理论分析结果表明，模型性能的退化程度与噪声强度呈正相关关系：在 MSE 损失下，收敛界随噪声方差增大而增大；在交叉熵损失下，收敛界随噪声率呈二次方增长。这一理论成果揭示了标签噪声导致 SFL 收敛震荡及准确率下降的根本原因，为后续鲁棒算法的设计提供了坚实的理论依据。
其次，在算法设计层面，本文提出了一种面向 Non-IID 数据与标签噪声协同治理的鲁棒 SFL 算法（SF-LCT）。为解决低质量数据导致的模型性能退化问题，该算法采用“检测 - 纠正 - 重训”的三阶段策略。在噪声检测阶段，设计了自适应机制，根据数据分布特性动态选择基于本地内在维度（LID）或标签熵（LE）的指标来精准定位含噪客户端；在噪声纠正阶段，利用干净客户端数据训练全局模型，进而修正含噪客户端的错误标签；最后在修正后的数据上进行传统 SFL 训练。在 CIFAR-10、CIFAR-100 及 Fashion-MNIST 等多个数据集上的实验结果表明，相较于基准算法，SF-LCT 将模型准确率最高提升了 20.01%。此外，基于树莓派搭建的真实物理测试床实测数据显示，该算法在保持高精度的同时，相较于现有去噪算法显著降低了能耗，验证了其在低质量数据环境下的鲁棒性与能效优势。
最后，在系统优化层面，本文彻底重构了 SFL 的训练范式，提出了资源受限场景下无反向传播的轻量化 SFL 框架（SOUL-SFL）。针对边缘设备算力弱、内存小及带宽受限的问题，该框架摒弃了客户端本地的反向传播计算，利用统一似然比（ULR）方法进行梯度估计，并将梯度估计的计算负载与显存需求完全卸载至服务器端。客户端仅需执行单次前向传播任务，从而将客户端的计算成本降低至原来的三分之一，并将激活值存储带来的内存占用降低至原来的二分之一。同时，优化了通信协议，将每轮次的数据传输量由“激活值 + 梯度”压缩为“激活值 + 模型参数”。实测数据显示，SOUL-SFL 在保持与标准 SFL 相当模型精度的前提下，计算成本、内存占用及通信开销最高分别降低了 59.38%、53.92% 和 58.97%。这一成果为在极端资源约束下的边缘设备部署高参数量模型提供了新的技术路径。
综上所述，本文通过理论分析、算法设计与系统构建三个层面的研究，形成了一套完整的面向低质量数据与资源受限场景的 SFL 协同优化方案，不仅丰富了分布式学习的理论体系，也为实际场景中的隐私保护机器学习应用提供了可行的解决方案。
\section{研究展望}
尽管本文在分割 - 联邦学习的收敛性分析、鲁棒性算法设计及轻量化框架构建方面取得了一定的研究成果，但受限于研究时间与个人能力，仍存在诸多值得进一步探索与完善的空间。针对当前研究中存在的局限性，结合分布式学习与边缘计算领域的最新发展趋势，未来的研究工作将重点围绕以下三个方面展开：
第一，深化模型收敛性理论分析，拓展损失函数适用范围并争取获得更紧的收敛界。本文目前的收敛性分析主要聚焦于均方误差与交叉熵这两种经典损失函数，虽然涵盖了回归与分类两大类任务，但面对日益复杂的深度学习应用场景，现有的损失函数体系尚显不足。未来计划针对更多类型的损失函数进行深入分析，例如针对类别不平衡场景的焦点损失（Focal Loss）、针对表征学习的对比损失（Contrastive Loss）以及针对生成任务的对抗损失等。不同的损失函数具有不同的几何性质与梯度特征，这将给收敛性证明带来新的挑战。此外，本文目前的收敛界推导过程中，为了满足理论分析的可解性，引入了一些较为宽松的假设条件（如梯度有界性、Lipschitz 连续性等），这在一定程度上导致了收敛界不够紧致。未来工作中，计划引入更先进的数学工具与非凸优化理论，尝试放宽这些假设条件，争取获得更接近实际训练轨迹的紧致收敛界，从而更精确地量化噪声与异构性对模型性能的影响边界，为算法超参数的选择提供更具指导意义的理论公式。
第二，增强低质量数据鲁棒性算法的适应性，设计能够应对更多种噪声模型的治理方案。本文提出的 SF-LCT 算法主要针对随机翻转噪声（Random Flipping Noise）进行了优化，并在 Non-IID 场景下取得了良好效果。然而，现实世界中的数据噪声形态远比随机翻转复杂多样。未来准备设计适应更多种噪声模型的鲁棒性算法，例如针对结构化噪声（Structured Noise）、对抗性噪声（Adversarial Noise）以及开放集噪声（Open-set Noise）的治理机制。结构化噪声往往与数据特征相关，对抗性噪声旨在恶意破坏模型训练，而开放集噪声则涉及未知类别的混入，这些场景下的噪声检测与纠正难度更大。未来的研究将探索结合元学习（Meta-Learning）与因果推断（Causal Inference）技术，使模型能够自适应地识别噪声类型并动态调整去噪策略。同时，考虑到隐私保护的限制，如何在密文域或联邦聚合过程中实现高效的噪声过滤，也是未来需要重点攻克的技术难点，旨在构建一个通用性强、适应性广的联邦噪声学习框架。
第三，进一步优化无反向传播 SFL 框架的资源效率，并在更广泛的模型与数据集上进行验证。本文提出的 SOUL-SFL 框架虽然通过服务器卸载与 ULR 方法显著降低了客户端负载，但仍存在优化空间。未来计划进一步优化该框架的计算、内存和通信负载。在计算层面，研究更低方差的梯度估计器，以减少 ULR 方法所需的噪声注入步数，从而降低服务器端的计算压力并间接减少通信频次；在内存层面，探索激活值压缩技术与动态模型切分策略，根据客户端实时资源状态自适应调整切分点，实现资源消耗与模型精度的最佳平衡；在通信层面，结合梯度稀疏化与量化技术，进一步压缩传输数据量。此外，目前的实验主要集中在 CNN 架构与图像分类数据集上。未来将在更多样的模型架构上进行实验，包括 Transformer 架构、图神经网络（GNN）以及新兴的大语言模型（LLM）微调场景。同时，拓展数据集的范围，涵盖医疗影像、物联网传感器数据及金融交易记录等具有高度隐私敏感性与资源约束特性的真实场景数据，以全面验证 SOUL-SFL 框架的通用性与实用性，推动其在工业互联网、智慧医疗及车联网等领域的实际落地应用。
总之，分割 - 联邦学习作为隐私保护与分布式计算相结合的前沿方向，其发展潜力巨大。希望通过未来的持续研究，能够进一步突破数据质量与系统资源的双重瓶颈，构建更加高效、稳健、安全的分布式智能学习系统，为人工智能技术的普惠化应用贡献力量。